% 第 4 章：换流器短路贡献
\section{换流器短路贡献}\label{sec:sc-converter}

\subsection{两类换流器短路模型}
换流器短路行为由 \fld{SCConverterModel}（\srcpath{short\_circuit.hpp:88}）区分：
\begin{itemize}
  \item \textbf{跟网（GridFollowing）}：受限电流源模型，IEC 60909 §6.7
        \begin{equation}
        I_{sc}=k\,I_{rated},
        \end{equation}
        限流倍数 $k$ 由换流器电流限 \fld{i\_limit\_pu} 给定，故障期间贡献被钳制于额定的 $k$ 倍；
  \item \textbf{构网（GridForming）}：内电势后接阻抗的电压源
        \begin{equation}
        I_{sc}=\frac{V_{int}}{Z_{filter}+Z_{virtual}},
        \end{equation}
        滤波阻抗与虚拟阻抗共同决定贡献。
\end{itemize}

\subsection{贡献结果}
逐换流器贡献记于 \fld{SCConverterContributionResult}（:144）：\fld{converter\_index}、\fld{bus\_id}、
\fld{model}（\texttt{grid\_following\_current\_source} / \texttt{ac\_grid\_forming\_voltage\_source}）、
\fld{ac\_grid\_forming}/\fld{dc\_grid\_forming} 标志、\fld{p\_rated\_mw}、\fld{i\_limit\_pu} 与
\fld{contribution\_ka}。该贡献并入故障母线 \fld{ikss\_converter\_contrib\_ka}（:114），与发电机/
电机/外网贡献分列，便于按来源审计。

\subsection{与旋转机的区别}
与同步机（内阻抗 $z_G''$ + 校正 $K_G$，见第~\ref{sec:sc-detailed}~章）不同，电力电子换流器的
短路电流受控制限流主导而非物理次暂态电抗，故跟网型贡献接近额定电流量级而非数倍。构网型
在暂态窗口内表现为受阻抗限制的电压源，其持续贡献取决于内电势与等效阻抗。二者均是稳态/
准稳态故障水平估计，\textbf{不}含控制环动态、FRT 轨迹与限流器切换时序。
